% Part: incompleteness
% Chapter: incompleteness-provability
% Section: second-incompleteness-thm

\documentclass[../../../include/open-logic-section]{subfiles}

\begin{document}

\olfileid[tr]{inc}{inp}{2in}

\olsection{İkinci Eksiklik Teoremi}

$\Th{PA}$'nın kendi tutarlılığını ispatlamadığı önermesini nasıl ifade
edebiliriz? $\Th{PA}$'nın tutarsız olduğunu söylemek,
$\Th{PA} \Proves \eq[0][1]$ demektir. Bu nedenle tutarlılık önermesi
$\OCon[\Th{PA}]$'yı $\lnot \OProv[\Th{PA}](\gn{\eq[0][1]})$ !!{sentence}
olarak alabiliriz; aşağıdaki teorem de istediğimiz sonucu verir:

\begin{thm}
\ollabel{thm:second-incompleteness} 
$\Th{PA}$ tutarlıysa $\Th{PA}$, $\OCon[\Th{PA}]$'yı !!{derive} etmez.
\end{thm}

Teoremin $\OCon[\Th{PA}]$'nın seçilen özel temsiline (yani
$\OProv[\Th{PA}](y)$'nin seçilen özel temsiline) bağlı olduğuna dikkat etmek
önemlidir. Yalnızca $\OProv[\Th{PA}](y)$ temsilinin üç !!{derivability}
koşulunu sağladığını kullanacağız; dolayısıyla teorem bu özelliklere sahip
!!a{derivability} yüklemi bulunan her kurama genellenir.

Teorem güzel bir can alıcı söz gibi ortaya çıktığından G\"odel'in argüman
taslağını okumak öğreticidir. Argüman şöyledir.
\olref[1in]{thm:first-incompleteness} ispatında kurduğumuz G\"odel tümcesi
$!G_\Th{PA}$ olsun. ``$\Th{PA}$ tutarlıysa $\Th{PA}$, $!G_\Th{PA}$'yı
!!{derive} etmez'' sonucunu göstermiştik. Bunu $\Th{PA}$ \emph{içinde}
biçimselleştirirsek
\[
\OCon[\Th{PA}] \lif \lnot \OProv[\Th{PA}](\gn{!G_\Th{PA}}).
\]
formülünün bir ispatını elde ederiz. Şimdi $\Th{PA}$'nın
$\OCon[\Th{PA}]$'yı !!{derive} ettiğini varsayın. O zaman
% [tr source emendation] Upstream switches here from the object-language
% provability formula \OProv to \Prov; the argument requires \OProv.
$\lnot \OProv[\Th{PA}](\gn{!G_\Th{PA}})$ formülünü de !!{derive} eder. Fakat
$!G_\Th{PA}$ bir G\"odel tümcesi olduğundan bu formül $!G_\Th{PA}$'ya
denktir. Öyleyse $\Th{PA}$, $!G_\Th{PA}$'yı !!{derive} eder.

Ama $\Th{PA}$ tutarlıysa $!G_\Th{PA}$'yı !!{derive} etmediğini biliyoruz!{}
Dolayısıyla $\Th{PA}$ tutarlıysa $\OCon[\Th{PA}]$'yı !!{derive} edemez.

Argümanı daha kesin kılmak için $!G_\Th{PA}$, $\Th{PA}$'nın G\"odel tümcesi
olsun. (P1)--(P3) !!{derivability} koşullarını kullanarak $\Th{PA}$'nın
$\OCon[\Th{PA}] \lif !G_\Th{PA}$ formülünü !!{derive} ettiğini göstereceğiz.
Bu, $\Th{PA}$'nın $\OCon[\Th{PA}]$'yı !!{derive} etmediğini gösterecektir.
Aşağıda ispatın $\Th{PA}$ içindeki bir taslağı verilmiştir. (Sadelik için
$\Th{PA}$ alt indislerini kaldırıyoruz.)
\begin{align}
& !G \liff \lnot \OProv(\gn{!G}) \ollabel{G2-1}\\
& \qquad\text{$!G$ bir G\"odel tümcesidir}\notag \\
& !G \lif \lnot \OProv(\gn{!G}) \ollabel{G2-2}\\
  & \qquad\text{\olref{G2-1}'den} \notag\\
& !G \lif
  (\OProv(\gn{!G}) \lif \lfalse) \ollabel{G2-3}\\
  & \qquad\text{\olref{G2-2}'den mantık aracılığıyla}\notag\\
& \OProv(\gn{
    !G \lif
    (\OProv(\gn{!G}) \lif \lfalse)
  }) \ollabel{G2-4}\\
  & \qquad\text{\olref{G2-3}'ten P1 koşulu aracılığıyla} \notag\\
& \OProv(\gn{!G}) \lif
  \OProv(\gn{
    (\OProv(\gn{!G}) \lif \lfalse)
    }) \ollabel{G2-5}\\
  & \qquad\text{\olref{G2-4}'ten P2 koşulu aracılığıyla} \notag\\
& \OProv(\gn{!G}) \lif (\OProv(\gn{\OProv(\gn{!G})}) \lif \OProv(\gn{\lfalse})) \ollabel{G2-6}\\
  & \qquad\text{\olref{G2-5}'ten P2 koşulu ve mantık aracılığıyla} \notag\\
& \OProv(\gn{!G}) \lif 
  \OProv(\gn{\OProv(\gn{!G})}) \ollabel{G2-7}\\
   & \qquad\text{P3 aracılığıyla} \notag\\
& \OProv(\gn{!G}) \lif \OProv(\gn{\lfalse}) \ollabel{G2-8}\\
  & \qquad \text{\olref{G2-6} ve \olref{G2-7}'den mantık aracılığıyla}\notag\\
& \OCon \lif \lnot \OProv(\gn{!G}) \ollabel{G2-9}\\
  & \qquad\text{\olref{G2-8}'in karşıt tersi ve $\OCon \ident \lnot \OProv(\gn{\lfalse})$}\notag \\
& \OCon \lif !G \notag\\
  & \qquad\text{\olref{G2-1} ve \olref{G2-9}'dan mantık aracılığıyla}\notag
\end{align}
Yukarıda mantığın kullanıldığı yerlerde yalnızca önerme mantığının temel
olguları kullanılmıştır. Örneğin \olref{G2-3}'te
$\Proves \lnot!A \liff (!A\lif \lfalse)$, \olref{G2-8}'de ise
$!A \lif (!B \lif !C), !A \lif !B \Proves !A \lif !C$ kullanılır.
\olref{G2-5} ve \olref{G2-6}'da P2 koşulu,
$\OProv(\gn{!A \lif !B}) \lif
(\OProv(\gn{!A}) \lif \OProv(\gn{!B}))$ biçimindeki P2 örneklerine dayanır.
Birincisinde $!A \ident !G$ ve
$!B \ident \OProv(\gn{!G}) \lif \lfalse$; ikincisinde ise
% [tr source emendation] Upstream drops the formula marker in \gn{G}; the
% second P2 instance requires the same \gn{!G} used throughout the proof.
$!A \ident \OProv(\gn{!G})$ ve $!B \ident \lfalse$ olur.

İkinci eksiklik teoreminin daha soyut biçimi şöyledir:

\begin{thm}
% [tr source emendation] Mere “axiomatized” is insufficient here; the proof
% relation must be computable.  The established token is axiomatizable.
\ollabel{thm:second-incompleteness-gen} $\Th{T}$, $\Th{Q}$'yu genişleten
herhangi bir tutarlı, !!{axiomatizable} kuram ve $\OProv[\Th{T}](y)$,
$\Th{T}$ için P1--P3 !!{derivability} koşullarını sağlayan herhangi bir formül
olsun. O zaman $\Th{T}$, $\OCon[\Th{T}]$'yi !!{derive} etmez.
% [tr source emendation] Upstream writes \OCon[T], dropping the theory macro.
\end{thm}

\begin{prob}
$\Th{PA}$'nın $!G_{\Th{PA}} \lif \OCon[\Th{PA}]$ formülünü !!{derive}
ettiğini gösterin.
\end{prob}

\begin{digress}
Çıkarılacak ders şudur: matematik için ``makul'' hiçbir tutarlı kuram kendi
tutarlılık önermesini !!{derive} edemez. $\Th{T}$, $\Th{Q}$'yu ve Hilbert'in
``sonlu yöntemlere dayanan'' akıl yürütmesini (bu her neyse) içeren bir
matematik kuramı olsun. O zaman $\Th{T}$'nin bütünü, $\Th{T}$'nin tutarlılık
önermesini !!{derive} edemez; dolayısıyla öncelikle onun sonlucu parçası da
$\Th{T}$'nin tutarlılık önermesini !!{derive} edemez. Bu anlamda ``bütün
matematik'' için sonlu yöntemlere dayanan bir tutarlılık ispatı bulunamaz.

``Sonlu yöntemlere dayanan'' teriminin yorumunda bir miktar esneklik vardır.
G\"odel, 1931 tarihli makalesinde, bizim böyle sayabileceğimiz bir şeyin
Hilbert'in biçimselleştirmek istediği matematik türlerinin dışında kalabileceği
olasılığını kabul eder. Fakat G\"odel cömert davranıyordu: bugün, makul biçimde
sonlu yöntemlere dayanan denebilecek fakat örneğin seçme aksiyomlu
Zermelo--Fraenkel küme kuramı $\Th{ZFC}$ içinde biçimselleştirilemeyecek bir
şeyi nasıl bulabileceğimizi görmek zordur.
\end{digress}

\end{document}
